\section{Dispersion Relations}
\label{sec:dispersion_relations}

A dispersion relation connects the angular frequency \(\omega\) of a mode with its wave vector \(\mathbf{k}\). It determines how individual phases move, how wave packets propagate, and whether different wavelengths travel at the same speed.

For the free real scalar field,

\[
\boxed{
\omega^{2}
=
c^{2}
\left(
|\mathbf{k}|^{2}+\mu^{2}
\right)
}.
\]

We focus on the positive-frequency branch

\[
\boxed{
\omega(k)
=
c\sqrt{k^{2}+\mu^{2}}
},
\]

where

\[
k=|\mathbf{k}|.
\]

\subsection{Massless dispersion}

When

\[
\mu=0,
\]

the relation becomes

\[
\boxed{
\omega=ck
}.
\]

This is linear in \(k\). Every Fourier component travels with the same characteristic speed \(c\), so an arbitrary one-dimensional profile may propagate without changing shape.

The graph of \(\omega(k)\) is a straight line through the origin.

\subsection{Massive dispersion}

For

\[
\mu>0,
\]

the relation is

\[
\boxed{
\omega=c\sqrt{k^{2}+\mu^{2}}
}.
\]

At

\[
k=0,
\]

the frequency is

\[
\boxed{
\omega_{0}=c\mu
}.
\]

Thus the graph begins at a nonzero frequency. The massive field has a frequency gap between the zero-frequency axis and the lowest positive-frequency mode.

\subsection{Phase velocity}

For a plane wave

\[
\cos(kx-\omega t),
\]

a surface of constant phase satisfies

\[
kx-\omega t=	ext{constant}.
\]

Differentiating gives

\[
k\frac{\mathrm{d}x}{\mathrm{d}t}-\omega=0.
\]

The phase velocity is therefore

\[
\boxed{
v_{\mathrm{ph}}
=
\frac{\omega}{k}
}.
\]

For a massless field,

\[
\boxed{
v_{\mathrm{ph}}=c}.
\]

For a massive field,

\[
\boxed{
v_{\mathrm{ph}}
=
c
\sqrt{1+
\frac{\mu^{2}}{k^{2}}}
}.
\]

Hence

\[
v_{\mathrm{ph}}>c
\]

for finite \(k\).

This does not imply superluminal signal propagation. A perfectly monochromatic plane wave extends throughout space and carries no localized front whose speed could transmit new information.

\subsection{Group velocity}

A localized wave packet is constructed from nearby wave numbers. Its envelope moves approximately with the group velocity

\[
\boxed{
v_{\mathrm{g}}
=
\frac{\mathrm{d}\omega}{\mathrm{d}k}
}.
\]

For the massive relation,

\[
\begin{aligned}
v_{\mathrm{g}}
&=
\frac{\mathrm{d}}{\mathrm{d}k}
\left[
c\sqrt{k^{2}+\mu^{2}}
\right]\\
&=
\frac{ck}{\sqrt{k^{2}+\mu^{2}}}.
\end{aligned}
\]

Therefore

\[
\boxed{
v_{\mathrm{g}}
=
\frac{c^{2}k}{\omega}
}.
\]

For \(\mu>0\),

\[
\boxed{
v_{\mathrm{g}}<c}.
\]

For the massless field,

\[
v_{\mathrm{g}}=c.
\]

\subsection{Relation between phase and group velocity}

For the free massive scalar field,

\[
\begin{aligned}
v_{\mathrm{ph}}v_{\mathrm{g}}
&=
\frac{\omega}{k}
\frac{c^{2}k}{\omega}\\
&=
c^{2}.
\end{aligned}
\]

Thus

\[
\boxed{
v_{\mathrm{ph}}v_{\mathrm{g}}=c^{2}
}.
\]

When the phase velocity exceeds \(c\), the group velocity is correspondingly smaller than \(c\).

\subsection{Connection with relativistic energy and momentum}

In a quantum interpretation, one associates

\[
E=\hbar\omega,
\qquad
\mathbf{p}=\hbar\mathbf{k}.
\]

Using

\[
\mu=rac{mc}{\hbar},
\]

the dispersion relation becomes

\[
E^{2}
=
|\mathbf{p}|^{2}c^{2}
+
m^{2}c^{4}.
\]

The group velocity then satisfies

\[
\boxed{
v_{\mathrm{g}}
=
\frac{c^{2}|\mathbf{p}|}{E}
}.
\]

For a relativistic particle,

\[
\mathbf{v}
=
\frac{c^{2}\mathbf{p}}{E}.
\]

Thus the group velocity of a wave packet agrees with the corresponding classical relativistic particle velocity. This is an interpretation of the same mathematical dispersion relation, not an assumption required for the classical field theory.

\subsection{Long-wavelength limit}

When

\[
k\ll\mu,
\]

write

\[
\omega
=
c\mu
\sqrt{1+
\frac{k^{2}}{\mu^{2}}}.
\]

Using

\[
\sqrt{1+z}
\approx
1+
\frac{z}{2}
\]

for small \(z\), we obtain

\[
\boxed{
\omega
\approx
c\mu
+
\frac{c k^{2}}{2\mu}
}.
\]

The frequency is dominated by the uniform oscillation frequency \(c\mu\). The group velocity is approximately

\[
\boxed{
v_{\mathrm{g}}
\approx
\frac{ck}{\mu}
}.
\]

Long-wavelength massive modes therefore move slowly.

\subsection{Short-wavelength limit}

When

\[
k\gg\mu,
\]

write

\[
\omega
=
ck
\sqrt{1+
\frac{\mu^{2}}{k^{2}}}.
\]

Then

\[
\boxed{
\omega
\approx
ck
+
\frac{c\mu^{2}}{2k}
}.
\]

The group velocity becomes

\[
\boxed{
v_{\mathrm{g}}
\approx
c
\left(
1-
\frac{\mu^{2}}{2k^{2}}
\right)
}.
\]

Short-wavelength massive modes approach massless propagation, but their group velocity remains slightly below \(c\).

\subsection{Wave-packet spreading}

A medium or equation is dispersive when different wave numbers have different group velocities. For the massive field,

\[
\frac{\mathrm{d}^{2}\omega}{\mathrm{d}k^{2}}
=
\frac{c\mu^{2}}
{(k^{2}+\mu^{2})^{3/2}}.
\]

This is nonzero for \(\mu>0\). Components with different \(k\) therefore move at different velocities, and a localized packet generally spreads as it evolves.

For the one-dimensional massless relation

\[
\omega=ck,
\]

the second derivative is zero. Ideal massless wave packets can propagate without dispersive spreading.

\subsection{A two-mode illustration}

Consider two nearby modes

\[
\phi
=
\cos(k_{1}x-\omega_{1}t)
+
\cos(k_{2}x-\omega_{2}t).
\]

Using the sum identity,

\[
\begin{aligned}
\phi
&=
2
\cos
\left[
\frac{(k_{1}-k_{2})x
-(\omega_{1}-\omega_{2})t}{2}
\right]\\
&\quad\times
\cos
\left[
\frac{(k_{1}+k_{2})x
-(\omega_{1}+\omega_{2})t}{2}
\right].
\end{aligned}
\]

The rapidly oscillating factor has phase velocity near

\[
\frac{\omega}{k},
\]

while the slowly varying envelope moves with

\[
\frac{\Delta\omega}{\Delta k}.
\]

In the limit of nearby modes,

\[
\frac{\Delta\omega}{\Delta k}
\longrightarrow
\frac{\mathrm{d}\omega}{\mathrm{d}k}
=
v_{\mathrm{g}}.
\]

\subsection{Causal propagation}

The group velocity is useful for narrow packets, but the strongest causal statement comes from the differential equation itself. The principal part is the d'Alembertian,

\[
\Box\phi,
\]

whose characteristics are the light cone. Compactly supported disturbances cannot create effects outside the future light cone when the initial-value problem is posed causally.

The lower-order mass term changes the interior evolution but not the characteristic cone.

\subsection{Negative-frequency branch}

The quadratic relation has two branches:

\[
\boxed{
\omega
=
\pm c\sqrt{k^{2}+\mu^{2}}
}.
\]

For a real classical field, positive- and negative-frequency complex modes are paired by complex conjugation. They are not independent real degrees of freedom.

In the trigonometric description, both branches are already contained in the sine and cosine time dependence.

\subsection{Frequency gap and static behaviour}

For a massive mode,

\[
|\omega|\geq c\mu.
\]

There are no oscillatory plane-wave modes with

\[
0<|\omega|<c\mu.
\]

If one instead fixes a low frequency below the gap and solves for a spatial wave number, then

\[
k^{2}
=
\frac{\omega^{2}}{c^{2}}
-
\mu^{2}<0.
\]

The wave number becomes imaginary, producing exponential spatial behaviour rather than propagation. This is the time-dependent analogue of the static decay length \(1/\mu\).

\subsection{Dimensional checks}

The quantities satisfy

\[
[k]=[\mu]=\mathrm{m}^{-1},
\]

and

\[
[\omega]=\mathrm{s}^{-1}.
\]

Thus

\[
c^{2}(k^{2}+\mu^{2})
\]

has dimensions \(\mathrm{s}^{-2}\), matching \(\omega^{2}\).

Both

\[
v_{\mathrm{ph}}=rac{\omega}{k}
\]

and

\[
v_{\mathrm{g}}=rac{\mathrm{d}\omega}{\mathrm{d}k}
\]

have dimensions of velocity.

\subsection{Common mistakes}

Typical errors include:

\begin{enumerate}
    \item identifying phase velocity with signal velocity;
    \item concluding that \(v_{\mathrm{ph}}>c\) violates relativity;
    \item differentiating \(\omega^{2}\) instead of \(\omega\) when computing group velocity;
    \item using the massless relation \(\omega=ck\) for a massive mode;
    \item forgetting the nonzero frequency at \(k=0\);
    \item claiming that the mass term changes the characteristic light cone;
    \item assuming that every dispersive wave packet propagates without changing shape.
\end{enumerate}

\subsection{What has been established}

The free scalar dispersion relation is

\[
\omega=c\sqrt{k^{2}+\mu^{2}}.
\]

Massless modes have equal phase and group velocities \(c\). Massive modes have phase velocity greater than \(c\), group velocity less than \(c\), and the product \(v_{\mathrm{ph}}v_{\mathrm{g}}=c^{2}\). The mass parameter creates a frequency gap and dispersive wave-packet evolution without changing the causal light cone. The next section derives the energy and momentum carried by scalar-field configurations.